% !PRINTED 219
\weberpage{219}
From this it follows that the division into rows is entirely independent
of the arbitrariness in the assumption about the
roots $\mathrm{sn}\,\Omega_1, \mathrm{sn}\,\Omega_2 \ldots$

\gesperrt{It follows further, however, that under the
substitutions of the group $\mathfrak{A}$ the rows are not torn apart,
but are only interchanged with one another.}
The group of the division equation is therefore imprimitive, and the
individual rows are the systems of imprimitivity (Bd.\,I, \S\,158).

For if one applies to $(\mu,\mu')$ and $(\nu,\nu')$ simultaneously the
substitution
\[
\begin{pmatrix} a, & b\\ c, & \partial \end{pmatrix}
\]
the congruence (3) remains preserved.

If one selects, among the substitutions of the group $\mathfrak{A}$, those substitutions
which interchange the roots of one row only among themselves, one obtains a group, and indeed a subgroup
of $\mathfrak{A}$. To every row there thus belongs such a subgroup of $\mathfrak{A}$.
If we denote by $R_{\mu,\mu'}$ the row in which the root $x_{\mu,\mu'}$ occurs,
and by $\mathfrak{A}_{\mu,\mu'}$ the subgroup of $\mathfrak{A}$ belonging to this row,
then by (3) [cf. \S\,63, (3)] the substitutions
\[
\begin{pmatrix} a, & b\\ c, & \partial \end{pmatrix}
\]
of $\mathfrak{A}_{\mu,\mu'}$ are characterized by the congruence

\weq{4}{(\partial\mu - b\mu')\mu' + (c\mu - a\mu')\mu \equiv 0 \pmod{n}}

In the same way one obtains a group $\mathfrak{B}_{\mu,\mu'}$ if one seeks out the substitutions
in $\mathfrak{B}$ which interchange the roots of the row $R_{\mu,\mu'}$ only
among themselves; their substitutions:
\[
\begin{pmatrix} \alpha, & \beta\\ \gamma, & \delta \end{pmatrix}
\]
are characterized by the two congruences:

\weq{5}{\begin{aligned}
(\delta\mu - \beta\mu')\mu' + (\gamma\mu - \alpha\mu')\mu &\equiv 0\\
\alpha\delta - \beta\gamma &\equiv 1
\end{aligned}\ \ (\mathrm{mod}\ n)}

For example, the groups $\mathfrak{A}_{1,0}$, $\mathfrak{B}_{1,0}$ consist of the substitutions
\[
\begin{pmatrix} a, & b\\ 0, & \partial \end{pmatrix}, \quad
\begin{pmatrix} \alpha, & \beta\\ 0, & \delta \end{pmatrix}, \quad \alpha\delta \equiv 1 \pmod{n}.
\]
